\begin{abstract}
This Master’s Thesis investigates the automated detection of self-disclosure in 16,000 chatbot interactions with elderly users. Grounded in Social Penetration Theory, the study evaluates various feature extraction pipelines using the Gemma-2-9b-it LLM. We compare binary versus rubric-based scoring and direct versus subtopic-based classification aggregated via Hadamard products. An unsupervised K-Means clustering approach is employed to identify thematic patterns, moving beyond heuristic metrics to a theory-driven $N+1$ architecture. Validation is conducted against a Gold Standard of 130 manually annotated sessions by four experts, utilizing Krippendorff’s Alpha, F1-score, and MCC. The results provide a robust methodological foundation for monitoring vulnerability and trust-building in conversational AI for senior citizens.

\vspace{0.5cm}
\noindent \textbf{Keywords:} Natural Language Processing, Self-Disclosure, LLM, Social Penetration Theory, K-Means Clustering, Elderly Care.
\end{abstract}
\newpage